\justifying
\NSFCBodyText

\subsubsection{已具备的实验条件}

申请人已确认可稳定使用 NVIDIA A40 显卡8张及 NVIDIA GeForce RTX 3090 Ti 显卡6张，并具备 Linux、Python、深度学习框架和 Git 等软件工具链。上述条件用于中小规模开源模型的参数高效后训练、轨迹生成与重复实验。不同型号显卡的显存和算力不混作同一计算单位，主要比较固定硬件类型并报告实际GPU时；公开模型能否满足预定批量和上下文长度，以首阶段显存测量为准。

研究数据以可公开获取、许可允许的文本工具任务及能源时序数据为基础。实验采用版本化、可重置的本地工具沙箱，保存任务状态、动作参数与判定结果，支持同任务重复运行和影子审计。申请人现有的时间序列与问答研究经验用于建立数据划分、版本管理和评价流程；专用于 PACT-Agent 的轨迹与适配器资产由本项目新建。

\subsubsection{尚缺少的实验条件及拟解决的途径}

当前需落实本项目人员投入与轨迹审计安排，具体为\textbf{【待补 F-PACT-02：人员名单、分工及投入期限】}。首阶段优先配置模型训练、工具环境和结果审计三类任务，按样本规模调整审计工作量；参与人数较少时减少扩展环境，保留主要机制对照。

需要依据正式基金附件核定模型生态适配、可公开数据范围与验收工作量，相关依据为\textbf{【待补 F-PACT-01：课题附件、经费与验收约定】}。公开基准依赖和模型许可先核验后使用，在线工具如出现不稳定或权限限制，切换到功能一致且可审计的本地版本，并将其列为定制实验而非原基准官方分数。验证器人工审计、公共模型下载和存储需求纳入首阶段资源核算。

\subsubsection{利用研究基地的计划与落实情况}

本项目核心实验依托申请人已确认可用的设备与软件条件开展。学院公共算力可作为补充申请渠道，尚未落实的固定账号、型号或额度不计入必需条件。项目不以新增国家实验室、全国重点实验室合作或私有行业数据授权为前提；如后续获得正式支持，在实际授权范围内开展扩展验证，并保留原有资源条件下的核心对照结果。
